\documentclass[11pt,a4paper]{article}

\usepackage{fontspec}
\usepackage[turkish]{babel}
\usepackage[a4paper,margin=2.15cm]{geometry}
\usepackage{setspace}
\usepackage{booktabs}
\usepackage{array}
\usepackage{tabularx}
\usepackage{caption}
\usepackage{float}
\usepackage[table]{xcolor}
\usepackage{tikz}
\usetikzlibrary{arrows.meta,positioning,shapes.geometric,fit}
\usepackage[colorlinks=true,linkcolor=black,urlcolor=blue,citecolor=black]{hyperref}

\setmainfont{Times New Roman}
\setsansfont{Arial}
\onehalfspacing

\newcommand{\code}[1]{\texttt{#1}}
\newcommand{\okcell}{\cellcolor{green!18}\textbf{Var}}
\newcommand{\partcell}{\cellcolor{yellow!25}\textbf{Kısmi}}
\newcommand{\nocell}{\cellcolor{red!9}Yok}
\newcommand{\metric}[2]{\textbf{#1}\\{\Large #2}}
\newcolumntype{Y}{>{\raggedright\arraybackslash}X}

\title{\textbf{Benzer Makaleler Karşılaştırma Raporu}\\
Saf LNN/CfC Mobil Robot Navigasyonu}
\author{Orkun Koray GÜNGÖR}
\date{11 Mayıs 2026}

\begin{document}
\shorthandoff{=}
\maketitle

\begin{abstract}
Bu raporda \code{benzer-makaleler} klasörüne indirilen dört çalışma incelendi ve
MuJoCo tabanlı saf LNN/CfC mobil robot navigasyonu çalışmasıyla karşılaştırıldı.
İncelenen literatür, LNN ailesinin robotik navigasyon için güçlü bir yöntemsel
temel sunduğunu gösteriyor; ancak okunan çalışmaların hiçbiri MuJoCo,
diferansiyel sürüşlü mobil robot, LiDAR/range gözlemi, saf CfC politika,
BC+DAgger eğitimi ve statik/dinamik harita değerlendirmesini aynı anda
birleştirmiyor. En yakın yöntemsel çalışma Vorbach ve ark. tarafından sunulan
görsel drone navigasyonu; en yakın platform/dinamik ortam çalışması ise Ao ve
ark. tarafından sunulan kalabalık ortamda NCP tabanlı mobil robot
navigasyonudur.
\end{abstract}

\section{Kapsam ve Kendi Çalışmamızın Profili}

Bu inceleme dört PDF ile sınırlıdır: Vorbach ve ark. (2021), Ao ve ark. (2026),
Tanna ve ark. (2024) ve Shang--Hu (2022). Amaç, bu çalışmaların bizim
\emph{saf LNN/CfC mobil robot navigasyonu} çalışmamıza hangi eksenlerde
benzediğini ve hangi eksenlerde ayrıldığını netleştirmektir.

\begin{table}[H]
\centering
\caption{Bu raporda referans alınan mevcut çalışma profili}
\begin{tabularx}{\linewidth}{@{}p{0.28\linewidth}Y@{}}
\toprule
\textbf{Öğe} & \textbf{Mevcut çalışma} \\
\midrule
Platform & MuJoCo ortamında diferansiyel sürüşlü mobil robot. \\
Gözlem & 38 boyutlu vektör: hedef göreli bilgileri, robot hızı, önceki aksiyon ve 32 ışınlı LiDAR/range taraması. \\
Politika & Saf LNN/CfC: \code{Linear+Tanh -> CfC -> CfC -> actor/critic}. \\
Eğitim & Behavior cloning + DAgger. Runtime sırasında waypoint takipçisi veya safety filter kullanılmıyor. \\
Değerlendirme & 22 custom statik harita + 2 dinamik harita. Son modelde 21/24 harita başarısı, 3/24 çarpışma, 0 timeout. \\
Model adı & \code{cfc\_deep192\_custom22\_dynamic\_dagger2}. \\
\bottomrule
\end{tabularx}
\end{table}

\section{Kısa Terminoloji Notu}

\textbf{LTC} bir sürekli zamanlı RNN hücresidir; gizli durumun değişimini ODE
benzeri bir dinamikle modeller ve zamansal bağımlılığı fiziksel sistemlere daha
yakın bir biçimde taşır. \textbf{CfC} aynı sürekli zamanlı fikri daha hızlı
çalışacak kapalı-form bir güncelleme ile uygular; bu nedenle pratik eğitim ve
inference maliyeti LTC'ye göre daha düşüktür. \textbf{NCP} ise tek başına ayrı
bir hücreden çok, biyolojik devrelerden esinlenen seyrek/yorumlanabilir bağlantı
topolojisi ve politika tasarımıdır; çoğu çalışmada NCP, LTC ya da CfC tabanlı
nöral dinamiklerle birlikte kullanılır.

Bu ayrım önemlidir: bazı makaleler ``NCP'' derken politika topolojisini, bazıları
``LTC'' derken ODE tabanlı hücreyi, bazıları da ``CfC'' derken kapalı-form
sürekli-zaman hücresini kasteder. Bizim son modelimiz NCP wiring kullanmadan,
saf ve derin CfC hücreleriyle kurulmuştur.

\section{İncelenen Makaleler}

\begin{table}[H]
\centering
\caption{Okunan dört makalenin hızlı özeti}
\small
\begin{tabularx}{\linewidth}{@{}p{0.06\linewidth}p{0.29\linewidth}p{0.21\linewidth}Y@{}}
\toprule
\textbf{No} & \textbf{Çalışma} & \textbf{Ana odak} & \textbf{Bizim işle ilişkisi} \\
\midrule
1 & Vorbach ve ark. (2021), \emph{Causal Navigation by Continuous-time Neural Networks} &
NCP/LTC/CfC ailesiyle görsel drone navigasyonu &
En güçlü yöntemsel akraba: sürekli zamanlı politika, imitation learning, kapalı çevrim navigasyon ve dinamik hedef takibi var; ancak platform drone ve sensör RGB kamera. \\
2 & Ao ve ark. (2026), \emph{Human-Centric Motion Planning in Crowded Spaces} &
SNCP-PPO ile kalabalık ortamda mobil robot navigasyonu &
En güçlü platform ve dinamik ortam benzeri: mobil robot, hareketli insan engelleri ve NCP tabanlı karar verme var; ancak eğitim PPO, temsil sosyal/state tabanlı ve yapı hibrit. \\
3 & Tanna ve ark. (2024), \emph{Optimizing Liquid Neural Networks} &
LTC ve CfC karşılaştırması, behavior cloning ve steering tahmini &
Mimari ve eğitim motivasyonu için yararlı: LTC/CfC karşılaştırması ve BC var; fakat robot navigasyonu veya LiDAR tabanlı kapalı çevrim değerlendirme yok. \\
4 & Shang ve Hu (2022), \emph{CALNet} &
LiDAR-kamera online kalibrasyonu için LTC içeren ağ &
LiDAR ile LTC'nin robotik algıda kullanıldığı yakın örnek; fakat politika, navigasyon ve kontrol problemi değil. \\
\bottomrule
\end{tabularx}
\end{table}

\section{Literatür Konumlandırma Haritası}

\begin{figure}[H]
\centering
\resizebox{\linewidth}{!}{%
\begin{tikzpicture}[
  x=1cm,y=1cm,
  point/.style={circle, draw=black!55, thick, minimum size=0.78cm, align=center, font=\sffamily\small},
  label/.style={align=left, font=\sffamily\small},
  axislabel/.style={font=\sffamily\small}
]
  \draw[->, thick] (0,0) -- (12.3,0) node[axislabel, below left=2pt and -8pt]
    {Platform/sensör benzerliği};
  \draw[->, thick] (0,0) -- (0,8.6) node[axislabel, above, rotate=90, anchor=south]
    {Politika/eğitim benzerliği};

  \node[axislabel, anchor=north west] at (0,-0.35) {Düşük};
  \node[axislabel, anchor=north east] at (12.1,-0.35) {Yüksek};
  \node[axislabel, anchor=east] at (-0.15,0.2) {Düşük};
  \node[axislabel, anchor=east] at (-0.15,8.15) {Yüksek};

  \draw[dashed, gray!55] (9.1,6.0) ellipse (2.3cm and 1.45cm);
  \node[label, text=gray!40!black] at (8.0,7.7) {Boşluk:\\mobile robot + LiDAR\\+ saf CfC + BC/DAgger};

  \node[point, fill=blue!18] (vorbach) at (2.3,7.2) {V};
  \node[label, right=0.12cm of vorbach] {Vorbach 2021\\drone + RGB + IL};

  \node[point, fill=green!18] (ao) at (8.6,4.4) {A};
  \node[label, right=0.12cm of ao] {Ao 2026\\mobile robot + crowd + PPO};

  \node[point, fill=orange!22] (tanna) at (4.4,5.3) {T};
  \node[label, right=0.12cm of tanna] {Tanna 2024\\LTC/CfC + BC benchmark};

  \node[point, fill=purple!17] (calnet) at (7.1,1.6) {C};
  \node[label, right=0.12cm of calnet] {CALNet 2022\\LiDAR + LTC calibration};

  \node[point, fill=red!20, minimum size=0.92cm] (ours) at (10.6,7.1) {\textbf{Biz}};
  \node[label, right=0.12cm of ours] {MuJoCo + diff-drive\\LiDAR + saf CfC\\BC+DAgger};
\end{tikzpicture}}
\caption{Dört makalenin mevcut çalışmaya göre konumu. Vorbach yönteme, Ao platforma,
Tanna mimari karşılaştırmaya, CALNet ise LiDAR/LTC sensör tarafına yakındır.}
\end{figure}

\clearpage
\section{Özellik Örtüşmesi}

\begin{table}[H]
\centering
\caption{Çalışmaların temel özelliklere göre karşılaştırılması}
{\scriptsize
\setlength{\tabcolsep}{2.5pt}
\renewcommand{\arraystretch}{1.15}
\begin{tabularx}{\linewidth}{@{}p{0.17\linewidth}p{0.09\linewidth}p{0.105\linewidth}p{0.085\linewidth}p{0.09\linewidth}p{0.10\linewidth}p{0.085\linewidth}Y@{}}
\toprule
\textbf{Çalışma} & \textbf{LNN/CT} & \textbf{Nav. politika} & \textbf{Mobil} & \textbf{LiDAR} & \textbf{Dinamik} & \textbf{IL/BC} & \textbf{Not} \\
\midrule
Bizim çalışma & \okcell & \okcell & \okcell & \okcell & \okcell & \okcell & Saf deep CfC, MuJoCo, 24 harita. \\
Vorbach 2021 & \okcell & \okcell & \nocell & \nocell & \okcell & \okcell & Drone ve RGB kamera; yöntemsel olarak çok yakın. \\
Ao 2026 & \okcell & \okcell & \okcell & \partcell & \okcell & \nocell & Sosyal/state temsili, PPO, SNCP ve graph/attention bileşenleri. \\
Tanna 2024 & \okcell & \partcell & \nocell & \nocell & \partcell & \okcell & Atari, steering ve GHI benchmark; kapalı çevrim robot testi yok. \\
CALNet 2022 & \okcell & \nocell & \partcell & \okcell & \partcell & \nocell & LiDAR-kamera kalibrasyonu; kontrol/navigasyon politikası değil. \\
\bottomrule
\end{tabularx}}
\end{table}

\section{Makale Bazlı Değerlendirme}

\subsection{Vorbach ve ark. (2021)}

Bu çalışma bizim işin en güçlü yöntemsel dayanağıdır. Yazarlar sürekli zamanlı
nöral ağların, özellikle NCP/LTC çizgisindeki modellerin, kapalı çevrim görsel
navigasyonda daha sağlam nedensel temsiller öğrenebildiğini savunur. Deneylerde
drone, RGB görüntülerden hedefe gitme, hareketli hedefi takip etme ve daha uzun
hafıza gerektiren görevlerde test edilir. Pasif validasyon kaybı iyi görünse bile
kapalı çevrim testte birçok modelin başarısız olması, bizim çalışmada da önemli
olan bir noktayı destekler: gerçek değerlendirme, modelin ortamla etkileşime
girdiği rollout üzerinde yapılmalıdır.

Bizim çalışmayla benzerlikler güçlüdür: sürekli zamanlı politika, imitation
learning, kapalı çevrim navigasyon ve dinamik nesne/engel karşılaşmaları. Temel
fark ise sensör ve platformdadır. Vorbach ve ark. drone + kamera ile çalışırken,
biz düşük boyutlu LiDAR/range vektörüyle diferansiyel sürüşlü mobil robotu test
ediyoruz. Bu nedenle bu makale ``aynı problem'' değil, ``aynı yöntem ailesinin
navigasyonda çalıştığını gösteren ana referans'' olarak kullanılmalıdır.

\subsection{Ao ve ark. (2026)}

Ao ve ark. makalesi platform açısından en yakın çalışmadır. Problem, kalabalık
insan ortamında mobil robotun güvenli ve sosyal olarak uyumlu hareket etmesidir.
SNCP-PPO yaklaşımı NCP tabanlı sürekli zamanlı bir ağ, actor-critic yapı, sosyal
baskı indeksi ve PPO eğitimini birleştirir. Simülasyon ve gerçek Turtlebot
deneylerinde, 10--20 kişilik dinamik senaryolarda yüzde 93--95 başarı raporlanır.

Bu çalışma bizim dinamik harita hedefimizle doğrudan konuşur: hareketli varlıklar,
çarpışma oranı, timeout ve başarı oranı aynı değerlendirme diline yakındır. Ancak
bizim çalışmadan önemli biçimde ayrılır. Ao ve ark. sosyal navigasyon problemine
odaklanır; gözlem temsili insan/robot state ve grafik etkileşimleri üzerinden
kurulur. Eğitim PPO ile yapılır ve model saf CfC değildir; NCP, attention,
spatio-temporal graph ve ödül tasarımı birlikte kullanılır. Bu nedenle Ao ve ark.
``mobil robot + dinamik ortam'' benzerliğini verir, fakat ``LiDAR + saf CfC +
BC/DAgger'' boşluğunu kapatmaz.

\subsection{Tanna ve ark. (2024)}

Tanna ve ark. çalışması doğrudan navigasyon makalesi değildir; fakat LTC ve CfC
hücrelerini pratik görevlerde karşılaştırdığı için mimari seçimimizi destekleyen
bir kaynaktır. Çalışmada Atari Breakout behavior cloning, otonom sürüş için
steering angle tahmini ve GHI tahmini incelenir. Sonuçlarda behavior cloning
görevinde LTC ve CfC'nin yakın performans verdiği, steering tahmininde CfC'nin
düşük MSE ile öne çıktığı, bazı zaman serisi görevlerinde ise LTC'nin daha iyi
olduğu raporlanır.

Bizim için ana çıkarım şudur: CfC seçimi yalnızca teorik bir tercih değildir;
BC/steering türü aksiyon tahmini görevlerinde de pratik olarak anlamlıdır. Ancak
Tanna ve ark. robotun kapalı çevrim ortamda çarpışma/timeout ürettiği bir
navigasyon problemi çözmez. Dolayısıyla bu makale yöntemsel destek ve
hiperparametre tartışması için kullanılmalı, mobil robot navigasyonunda doğrudan
rakip çalışma gibi sunulmamalıdır.

\subsection{Shang ve Hu (2022) -- CALNet}

CALNet, 3D LiDAR ve 2D kamera arasındaki online ekstrinsik kalibrasyonu öğrenen
bir ağdır. Yapıda channel attention, hybrid spatial pyramid pooling ve LTC
bileşeni birlikte kullanılır. Çalışma, LiDAR tabanlı robotik algı ve LTC'nin aynı
sistemde yer aldığı önemli bir örnektir; ayrıca robotik algı tarafında LNN
ailesinin kullanılabileceğini gösterir.

Fakat CALNet bizim navigasyon çalışmamıza yalnızca sensör/mimari kesişiminden
yakındır. Ağ bir politika üretmez, robot kontrolü yapmaz ve başarı/çarpışma gibi
navigasyon metrikleriyle değerlendirilmez. Bu yüzden CALNet, ``LiDAR + LTC''
bağlantısını göstermek için faydalıdır; ama bizim ana katkımız olan saf CfC
politika ile LiDAR tabanlı mobil robot navigasyonu bu çalışmada yoktur.

\section{Benzerlik Puanlaması}

\begin{table}[H]
\centering
\caption{Nitel benzerlik puanlaması: 1 düşük, 5 yüksek}
\small
\begin{tabularx}{\linewidth}{@{}p{0.18\linewidth}ccccY@{}}
\toprule
\textbf{Çalışma} & \textbf{Yöntem} & \textbf{Platform} & \textbf{Sensör} & \textbf{Eğitim} & \textbf{Yorum} \\
\midrule
Vorbach 2021 & 5 & 1 & 1 & 5 & Saf/kompakt sürekli zamanlı politika ve IL nedeniyle ana yöntem referansı. \\
Ao 2026 & 4 & 5 & 2 & 1 & Mobil robot ve dinamik insan ortamı çok yakın; PPO ve sosyal graph yapısı farklı. \\
Tanna 2024 & 4 & 1 & 1 & 4 & LTC/CfC ve BC deneyleri yararlı; robot navigasyonu yok. \\
CALNet 2022 & 2 & 2 & 5 & 1 & LiDAR + LTC kesişimi var; kalibrasyon problemi, kontrol değil. \\
\bottomrule
\end{tabularx}
\end{table}

\section{Net Bulgular}

\begin{enumerate}
  \item \textbf{Doğrudan eşleşen çalışma bulunmadı.} Okunan dört makalenin hiçbiri
  MuJoCo + diferansiyel sürüşlü mobil robot + LiDAR/range gözlemi + saf CfC
  politika + BC/DAgger + dinamik harita değerlendirmesini aynı anda kapsamıyor.

  \item \textbf{Yöntemsel çizgi güçlü.} Vorbach ve ark. sürekli zamanlı/nöral
  circuit politikaların kapalı çevrim navigasyonda neden anlamlı olduğunu
  gösteriyor. Bu, bizim saf CfC politikasını savunurken kullanılacak en güçlü
  temel referanstır.

  \item \textbf{Dinamik mobil robot karşılığı var ama hibrit.} Ao ve ark. mobil
  robot + dinamik insan ortamı için çok yakın bir problem çözüyor; fakat model
  saf CfC değil ve eğitim imitation learning değil. Bu fark, bizim çalışmanın
  daha sade ve izole edilmiş LNN/CfC katkısını belirginleştirir.

  \item \textbf{LiDAR tarafında boşluk devam ediyor.} CALNet LiDAR ile LTC'yi
  birleştirir; ancak bu bir kalibrasyon ağıdır. LiDAR/range vektöründen doğrudan
  hız komutu üreten saf CfC mobil robot politikası literatürde bu dört çalışma
  içinde görünmüyor.

  \item \textbf{BC+DAgger konumlandırması makul.} Tanna ve ark. BC görevlerinde
  LTC/CfC'nin kullanılabilirliğini destekler; Vorbach ve ark. ise imitation
  learning ile navigasyonu destekler. Bizim DAgger eklememiz, kapalı çevrim
  dağılım kaymasına karşı pratik bir güçlendirme olarak konumlandırılabilir.
\end{enumerate}

\section{Tez/Makale İçin Kullanılabilecek Konumlandırma}

Bu çalışma, sürekli zamanlı nöral politikaların navigasyon için etkili
olabileceğini gösteren drone + RGB literatüründen ayrılarak, düşük boyutlu
LiDAR/range gözlemiyle diferansiyel sürüşlü mobil robot navigasyonuna odaklanır.
Mevcut yakın çalışmalar ya görsel drone navigasyonu, ya sosyal mobil robot
navigasyonu, ya LTC/CfC benchmarkları, ya da LiDAR-kamera kalibrasyonu
çerçevesindedir. Bu nedenle çalışmanın katkısı, saf CfC politikanın MuJoCo
ortamında statik ve dinamik haritalar üzerinde BC+DAgger ile eğitilip
değerlendirilmesi olarak konumlandırılmalıdır.

\begin{table}[H]
\centering
\caption{Çalışmanın literatürdeki boşluğa cevabı}
\begin{tabularx}{\linewidth}{@{}p{0.24\linewidth}Y@{}}
\toprule
\textbf{Literatürde görülen durum} & \textbf{Bizim çalışmanın cevabı} \\
\midrule
LNN navigasyon çoğunlukla drone + RGB üzerinden gösterilmiş. &
Mobil robot ve LiDAR/range vektörüyle aynı sürekli zamanlı politika ailesi test ediliyor. \\
Mobil robot dinamik ortam çalışmaları genellikle sosyal/RL/graph tabanlı. &
Saf CfC politika izole edilerek, ek graph/attention/safety filter olmadan deneniyor. \\
LiDAR + LTC örnekleri algı/kalibrasyon tarafında kalıyor. &
LiDAR doğrudan kontrol politikasının gözlemi olarak kullanılıyor. \\
Kapalı çevrim dağılım kayması çoğu BC çalışmasında sorun. &
BC sonrası DAgger ile modelin kendi rollout dağılımından veri toplanıyor. \\
\bottomrule
\end{tabularx}
\end{table}

\section{Kaynakça}

\begin{enumerate}
  \item Ao, T., Li, H., Tian, Y., Fu, L., Leng, Y., Shi, H., Shi, L., \& Zhou, Y. (2026).
  Human-Centric Motion Planning in Crowded Spaces: A Structured Neural Circuit
  Approach with Social Interaction-Awareness. \emph{International Journal of
  Social Robotics}, 18, 52. \url{https://doi.org/10.1007/s12369-026-01389-9}

  \item Shang, H.-C., \& Hu, B.-J. (2022). CALNet: LiDAR-Camera Online
  Calibration With Channel Attention and Liquid Time-Constant Network. In
  \emph{2022 26th International Conference on Pattern Recognition (ICPR)},
  5147--5154. \url{https://doi.org/10.1109/ICPR56361.2022.9956145}

  \item Tanna, P. D., Khare, S., Jain, S. S., Vishnoi, S. R., Roy, S., \& Challa, J. S.
  (2024). Optimizing Liquid Neural Networks: A comparative study of LTCs and CFCs.
  In \emph{2024 IEEE International Conference on Big Data (BigData)}, 1470--1475.
  \url{https://doi.org/10.1109/BigData62323.2024.10826128}

  \item Vorbach, C., Hasani, R., Amini, A., Lechner, M., \& Rus, D. (2021).
  Causal Navigation by Continuous-time Neural Networks. In \emph{Advances in
  Neural Information Processing Systems 34}.
\end{enumerate}

\end{document}
